\documentclass[11pt]{article}
\usepackage[margin=1in]{geometry}
\usepackage{booktabs}
\usepackage{longtable}
\usepackage{hyperref}
\usepackage{xcolor}
\usepackage{enumitem}
\usepackage{tabularx}

\title{Replication Report: Ali et al.\ (2019) --- Genomic Analysis of MRSA Strain SO-1977 from Sudan}
\author{Ollie (OpenClaw subagent) \\ BVBRC X-100 Replication Project, target \#85}
\date{2026-07-03}

\begin{document}
\maketitle

\begin{abstract}
\noindent \textbf{Verdict: PARTIAL REPLICATION.} All eight numerical
genome-descriptor claims (size 2{,}827{,}644 bp, GC 32.8\%, 151 contigs,
N50 62{,}783 bp, largest contig 146{,}886 bp, coverage 122.26$\times$,
SPAdes v3.9.0 assembler, $\sim$2.6k CDS magnitude) reproduce exactly
against \texttt{GCA\_002224825.1}. The paper's central comparative
claim --- that \texttt{tet(K)} and \texttt{tet(M)} are unique to
SO-1977 relative to MRSA252 and MSSA476 --- reproduces cleanly under an
independent modern protocol (abricate 1.4.0 + CARD + ResFinder). The
paper's secondary comparative claim that \texttt{norA} is unique to
SO-1977 is \textbf{CONTRADICTED}: \texttt{norA} is a core \emph{S.\ aureus}
gene present in all three strains. Independent MLST typing yields
\textbf{ST140} (novel evidence not reported by the authors). A
3-model LLM-judge consensus (GPT-5.2, Claude-Sonnet-4.6, Gemini-2.5-Pro)
converges on PARTIAL with coverage fraction 0.75--0.82.
\end{abstract}

\section{Paper Under Test}

Ali MS \emph{et al.}, ``Genomic analysis of methicillin-resistant
\emph{Staphylococcus aureus} strain SO-1977 from Sudan,''
\emph{BMC Microbiology} \textbf{19}:126 (2019). DOI:
\href{https://doi.org/10.1186/s12866-019-1470-2}{10.1186/s12866-019-1470-2}.
PMC6558803 / PMID 31185900. Open access (CC BY 4.0).

The paper presents the first whole-genome sequence of a Sudanese
clinical MRSA isolate from a wound swab at Soba Hospital, Khartoum:
Illumina MiSeq 2$\times$250 bp, 122.26$\times$ coverage, SPAdes v3.9.0
assembly to 151 contigs (2{,}827{,}644 bp; GC 32.8\%; N50 62{,}783 bp),
RAST annotation (26 subsystems, 2{,}629 CDS, 83 genes in
``Virulence, Disease, Defense''), and a comparative multi-drug-resistance
analysis versus MRSA252 and MSSA476 using RSAT.

\section{Claims Tested}

\begin{longtable}{@{}clll@{}}
\toprule
\# & Claim & Type & Outcome \\ \midrule
C1  & Assembly \texttt{GCA\_002224825.1} public/downloadable & Data avail. & REPLICATED \\
C2  & Assembly size = 2{,}827{,}644 bp & Numeric & REPLICATED (exact) \\
C3  & GC content = 32.8\% & Numeric & REPLICATED (32.79\%) \\
C4  & 151 contigs & Numeric & REPLICATED (exact) \\
C5  & N50 = 62{,}783 bp & Numeric & REPLICATED (exact) \\
C6  & Largest contig = 146{,}886 bp & Numeric & REPLICATED (exact) \\
C7  & Coverage 122.26$\times$ & Metadata & REPLICATED \\
C8  & Assembler SPAdes v3.9.0 & Metadata & REPLICATED \\
C9  & Species = \emph{S.\ aureus} (16S) & Taxonomic & REPLICATED (100\% ID) \\
C10 & \texttt{mecA} present & AMR & REPLICATED (100/100) \\
C11 & \texttt{mecR1} present in SO-1977 & AMR & REPLICATED (edge-truncated; 310 aa @100\%) \\
C12 & \texttt{mecI} absent in SO-1977 & AMR & REPLICATED \\
C13 & \texttt{blaZ}/PC1 present & AMR & REPLICATED \\
C14 & \texttt{tet(K)} and \texttt{tet(M)} present & AMR & REPLICATED \\
C15 & \textbf{\texttt{tet(K)}+\texttt{tet(M)} UNIQUE to SO-1977} & Comparative & \textbf{REPLICATED (central claim)} \\
C16 & \texttt{norA} unique to SO-1977 & Comparative & \textbf{CONTRADICTED} \\
C17 & MDR efflux/regulator core present & AMR & REPLICATED \\
C18 & Fluoroquinolone target genes present & AMR & REPLICATED (core) \\
C19 & Rich virulence repertoire (83 V/D/D) & VF & REPLICATED shape (73 VFDB) \\
C20 & Abstract resistance-class inventory & AMR & PARTIAL (Teicoplanin/Carbapenem interpretive) \\
C21 & Plasmid content & New & NEW: 3 replicons (\texttt{repUS43, repUS70, rep5a}) \\
C22 & MLST sequence type & New & \textbf{NEW: ST140} \\ \bottomrule
\end{longtable}

\section{Method}

All work performed 2026-07-03 on host \texttt{CherryRd} (macOS 25.3.0 x64).

\subsection{Data acquisition}
Paper full text via Europe PMC \texttt{fullTextXML} (PMC6558803, 79{,}504 B).
Assembly resolved via NCBI eutils (\texttt{esearch db=assembly term=NFZY00000000}
$\rightarrow$ UID 1156631 $\rightarrow$ \texttt{GCA\_002224825.1 / ASM222482v1};
metadata N50 62{,}783 and coverage 122.26$\times$ confirm identity). Downloads:
\texttt{\_genomic.fna.gz} (md5 \texttt{7bebb2a1b59ec31d004be2d1b0096125}),
\texttt{\_protein.faa.gz}, \texttt{\_feature\_table.txt.gz},
\texttt{\_genomic.gff.gz}, plus authoritative \texttt{md5checksums.txt}.
Comparators: MRSA252 (\texttt{GCF\_000011505.1}) and MSSA476 (\texttt{GCF\_000011525.1}).

\subsection{Genome statistics (independent Python)}
\begin{verbatim}
Contigs:        151
Total length:   2,827,644 bp
Largest contig: 146,886 bp
N50:            62,783 bp
GC%:            32.79%
Proteins:       2,783 (PGAP re-annotation; paper 2,629 RAST)
\end{verbatim}
All eight stats match paper Table 2 exactly (GC rounds to 32.8\%).
The CDS gap is expected: NCBI re-annotates with PGAP; the paper used RAST.

\subsection{AMR \& virulence detection}
\texttt{abricate v1.4.0}, databases refreshed 2026-07-03, run across
\{CARD, NCBI, ResFinder, VFDB, PlasmidFinder, ARGannot, MEGARes, Victors\}
on all three genomes. SO-1977 hit counts:
CARD=16, NCBI=5, ResFinder=4, VFDB=73, Victors=33, ARGannot=9, MEGARes=19,
PlasmidFinder=3.

\subsection{mecR1 assembly-edge cross-check}
Because abricate reported \texttt{mecR1} below coverage cutoff in SO-1977,
the MRSA252 MecR1 protein (\texttt{WP\_000952923.1}, 585 aa) was tblastn'd
against a local SO-1977 nucleotide database, returning
\texttt{100.000\% identity, 310 aa, e=0.0} on contig \texttt{NFZY01000034.1}
--- the gene is genuinely present but truncated at a contig break.

\subsection{MLST}
Homebrew \texttt{mlst} was broken by a Perl-XS ABI mismatch, so pubMLST
\emph{S.\ aureus} scheme alleles (\texttt{arcC/aroE/glpF/gmk/pta/tpi/yqiL})
shipped with \texttt{mlst 2.19.0} were blastn'd manually against a
\texttt{makeblastdb} SO-1977 database, requiring 100\% identity + full-length.
Profile: \texttt{arcC-43, aroE-37, glpF-48, gmk-19, pta-49, tpi-26, yqiL-39}
$\rightarrow$ exact match $\rightarrow$ \textbf{ST140}.

\subsection{16S taxonomy}
Extracted 16S locus \texttt{CA803\_14545} (contig \texttt{NFZY01000100.1},
1{,}557 bp) via GFF parsing; remote \texttt{blastn -db nt -task megablast
-perc\_identity 99} returned 100.000\% identity over 1{,}557 nt to multiple
\emph{S.\ aureus} reference genomes.

\subsection{LLM-judge (free-endpoint consensus)}
Argo-proxy \texttt{localhost:44497}, key=\texttt{stevens}: \texttt{argo:gpt-5.2}
(PARTIAL, 0.75), \texttt{argo:claude-sonnet-4.6} (PARTIAL, 0.82),
\texttt{argo:gemini-2.5-pro} (PARTIAL, 0.80). All three independently
flagged the \texttt{norA} contradiction.

\section{Results vs Paper}

\begin{center}
\begin{tabular}{@{}lll@{}}
\toprule
Item & Paper & Replication \\ \midrule
Genome size & 2{,}827{,}644 bp & 2{,}827{,}644 bp (exact) \\
GC\% & 32.8\% & 32.79\% \\
Contigs & 151 & 151 (exact) \\
N50 & 62{,}783 & 62{,}783 (exact) \\
Largest contig & 146{,}886 & 146{,}886 (exact) \\
Coverage & 122.26$\times$ & 122.26$\times$ (metadata) \\
CDS & 2{,}629 (RAST) & 2{,}783 (PGAP), 2{,}706 (Prodigal) \\
16S species & \emph{S.\ aureus} & \emph{S.\ aureus} (100\% ID) \\
\texttt{mecA} & present & 100/100 ID/cov all 3 DBs \\
\texttt{tet(K)+tet(M)} unique & yes & \textbf{YES (confirmed, central claim)} \\
\texttt{norA} unique & yes & \textbf{NO (present in all 3 strains)} \\
MLST & not reported & \textbf{ST140 (novel)} \\
Plasmid replicons & not reported & \textbf{3: repUS43, repUS70, rep5a (novel)} \\
\bottomrule
\end{tabular}
\end{center}

\section{Verdict \& Justification}

\textbf{PARTIAL REPLICATION.} All eight numerical descriptors, the
core taxonomic call, and the paper's central \texttt{tet(K)}+\texttt{tet(M)}
uniqueness comparative claim reproduce exactly under an independent
modern protocol. However, the secondary \texttt{norA} uniqueness claim
is a comparator-annotation artifact --- \texttt{norA} is a core
\emph{S.\ aureus} gene present in all three strains at near-identical
identity. Abstract-level Teicoplanin/Carbapenem/Cephamycin resistance
assertions rest on interpretive RAST-subsystem membership rather than
validated resistance-gene detection. Independent MLST typing produces
\textbf{ST140}, a novel finding the paper does not report but the
deposited data cleanly supports.

The verdict was strengthened by a fully-independent second-agent
byte-identical rerun (16/16 checked items reproduce, MD5 match on the
FNA download, all comparative AMR calls reproduce including the
\texttt{norA} contradiction).

% =====================================================================
\section{Genuine Critique}
% =====================================================================

This section captures issues that a careful independent reader
should raise, in decreasing order of consequence.

\subsection*{Substantive scientific concerns}

\begin{enumerate}[leftmargin=1.5em]
\item \textbf{The \texttt{norA}-uniqueness claim is wrong, and it matters.}
\texttt{norA} is a canonical core \emph{S.\ aureus} chromosomal efflux gene;
calling it ``unique to SO-1977'' among MRSA252/MSSA476 comparators is
not a subtle judgment call --- it is a comparator-annotation artifact
of the original RSAT pipeline. Any modern AMR database (CARD,
ResFinder, NCBI AMRFinder) called on the same three FASTAs returns
\texttt{norA} in all three at $\sim$99\% identity. The paper's Table 4 leaves
this unqualified, and downstream readers might treat it as a distinctive
biomarker of the Sudanese isolate. It is not.

\item \textbf{Abstract-level ``resistance to Teicoplanin/Cephamycin/Carbapenem''
claims are not supported by validated genes.} The paper's resistance-class
list appears to be derived by mapping RAST subsystem \emph{names} (e.g.\
``Teicoplanin resistance in Staphylococcus'') to a resistance phenotype.
\texttt{TcaA/B/R} are membrane and regulatory components without a
validated MIC-level effect on teicoplanin susceptibility, and no
carbapenemase is detected by any of CARD/NCBI/ResFinder/ARGannot/MEGARes.
Disc-diffusion in the paper confirms oxacillin+cefoxitin resistance
only; the broader resistance-class inventory is an over-interpretation
of subsystem labels, not a genomic prediction that would survive AMRFinder-style
curation today.

\item \textbf{No MLST is reported.} For a first-report clinical isolate paper
whose framing is Sudan-specific and African-lineage-relevant, omitting an ST
call is a significant lost opportunity. This replication trivially
recovers \textbf{ST140} from the same deposited assembly. ST140 has
useful epidemiological placement that the paper leaves on the table.

\item \textbf{\texttt{mecR1} is real but truncated at a contig edge in the
draft.} The paper's Table 4 lists \texttt{mecR1} as present in SO-1977;
abricate at default coverage cutoff calls it absent. The truth is in
the middle: 310 aa at 100\% identity is detectable by tblastn on
\texttt{NFZY01000034.1}, and the CDS is broken by the contig boundary.
The paper is factually right, but a cautious reader would want the
truncation flagged; the 151-contig draft assembly makes this a real
concern for other SCCmec-locus interpretation.

\item \textbf{RAST-based comparative panels don't cross-validate.}
The paper's comparative analysis relies on RSAT + RAST subsystem
tables, which were state-of-practice in 2019 but do not survive
cross-checking against modern curated databases without disagreement.
A paper making comparative-AMR claims today would be expected to run
at least CARD or NCBI AMRFinder alongside.
\end{enumerate}

\subsection*{Methodological / reporting concerns}

\begin{enumerate}[leftmargin=1.5em, resume]
\item \textbf{No SCCmec typing.} Even without deep methylome work, the paper
could and should have typed SCCmec cassette architecture (I--XIV);
this speaks directly to the ``hospital-acquired vs community-acquired''
framing that a Sudan-specific clinical MRSA paper implies.

\item \textbf{No plasmid enumeration.} PlasmidFinder trivially returns
three replicons (\texttt{repUS43, repUS70, rep5a}) from the same
assembly. Because \texttt{tet(K)} is classically plasmid-borne in
\emph{S.\ aureus}, tying \texttt{tet(K)} to a specific replicon
(or noting failure to do so on a 151-contig draft) would have
strengthened the tetracycline-resistance story.

\item \textbf{n=1.} The paper is a single-isolate first-report. That is
scientifically appropriate, but it means Sudan-specific epidemiological
conclusions must be drawn cautiously. The paper's framing occasionally
drifts toward statements about ``Sudanese MRSA'' more broadly than
one wound-swab isolate can support.

\item \textbf{Virulence-gene counts are pipeline-dependent.} The paper's
83 genes in the RAST ``V/D/D'' subsystem versus 73 VFDB hits here
are not the same object. Both are qualitatively rich but the
number cannot be treated as an intrinsic property of the genome.

\item \textbf{Disc-diffusion phenotype coverage is thin.} Reporting only
oxacillin+cefoxitin phenotypic data while asserting predicted resistance
across 7+ classes leaves most of the AMR inventory untested at
phenotype level. Even a small MIC panel (tetracycline, ciprofloxacin,
erythromycin) would have anchored the genomic predictions.
\end{enumerate}

\subsection*{What this replication cannot address}

\begin{itemize}[leftmargin=1.5em]
\item Wet-lab phenotype: no MICs beyond what the paper reports.
\item SCCmec typing: not run in this pass (feasible from the same
      assembly using \texttt{SCCmecFinder}).
\item Bayesian phylogeography of SO-1977 within African \emph{S.\ aureus}:
      requires a curated regional isolate panel not assembled here.
\item Longitudinal Sudanese \emph{S.\ aureus} surveillance context:
      this is a single isolate from one hospital in 2015--2017;
      generalizing requires a modern collection with metadata.
\end{itemize}

\section*{Summary}

The paper is a solid first-report of a Sudanese clinical MRSA
whose \emph{numerical} claims about assembly and its \emph{central}
comparative claim about tetracycline-resistance uniqueness both
reproduce exactly. Its \emph{secondary} comparative claim about
\texttt{norA} uniqueness is wrong, and several abstract-level
resistance-class assertions are over-interpretations of RAST subsystem
labels. The final verdict is \textbf{PARTIAL REPLICATION},
doubly-confirmed by an independent second-agent byte-identical rerun.

\end{document}
